\scanpage{65}
% archive sheet number: 64

% L1–2, heading and form
Faisceaux $l$-élémentaires sur $X$ [$X$ noeth]

% L2–7, definition
De la forme $i_*(G)_{U,X}$, où

% L3, canonical immersion
\[
i:\operatorname{Spec} k(x)\longrightarrow X\qquad\text{imm. canonique}\quad(x\in X)
\]

% L4–7, sheaf and open
$G$ un faisceau fini de $l$-torsion sur $X$

$U$ un \inserted{voisinage} ouvert de $x$ dans $Y=\overline{x}$, $U=X-Z$ \deleted{Faisceau strictement élémentaire : \unsure{on}} tels que $G$ est \uline{étale} sur $U$, et $U$ est \uline{unibranche}.

% L8–10, proposition
\uline{Proposition} Tout faisceau de $l$-torsion noeth sur $X$ a une suite de composition formée de faisceaux de $l$-torsion élémentaires stricts.

% L11, corollary heading
\uline{Cor 1} Soit $n$ un entier. Conditions équivalentes

% L12, first condition
(i) $\mathrm{cd}_l(X)\leq n$

% L13–16, second condition
(ii) Pour tt $Y$ intègre fermé dans $X$, \deleted{\unsure{et}} \inserted{pt gén $x$,} tt fermé $Z$ de $Y$, \deleted{\unsure{d'où|faisceau}} tout $k(x)$-faisceau $G$ fini de $l$-torsion, si $i$ désigne l'immersion canonique $\operatorname{Spec} k(x)\longrightarrow Y$, on a

% L17–18, tests a and b
\[
\begin{aligned}
\text{a)}\quad&H^n(Y,i_*(G))\longrightarrow H^n(Z,i_*(G)|Z)\quad\text{surjectif}\\
\text{b)}\quad&H^{n+1}(Y,i_*(G))=0
\end{aligned}
\]

% L19–22, finite variant
\deleted{De plus, on peut} (ii bis) \unsure{comme} ii), sauf que $Y\longrightarrow X$ est fini, \unsure{génériqt|finissant} une immersion, $Y$ intègre \uline{unibranche}, $i_*(G)$ \uline{non ramifié} sur $Y-Z$.

% L23–25, cancelled corollary
\deleted{\uline{Cor. 1 bis} Conditions équivalentes (i) Pour tt $F$ de $l$-torsion, et $T$ fermé, $H^i_T(X,F)=0$ si $i>n$ (ii) \unsure{idem}}

% L24–25, cancelled support formulas
\[
\left\{\begin{aligned}
&\deleted{\text{a)}\ H^n_T(Y,i_{\unsure{*|d}}(G))\to H^n_{T\cap Z}(Z,i_*(G)|Z)\quad\text{surj}}\\
&\deleted{\text{b)}\ H^{n+1}_T(Y,i_*(G))=0}
\end{aligned}\right.
\]

% Foot, left margin
% slanted
\margin{Se donner une famille de $T$, contenant $X$, stable par \uline{\unsure{trans|tr. invs|unions}}}
